\label{sec:author_evaluation}
We conducted an internal review of the generated papers with authors.
Recent AI-generated papers include some degree of manual post-editing~\citep{zochi2025, weng2025deepscientist}, and only a few studies have carefully examined the raw, unedited outputs of AI systems~\citep{yamada2025ai}.
However, evaluating the raw, unedited outputs is essential for accurately understanding the current limitations of AI Scientist systems.

The review here does not evaluate whether the paper has the level of contribution, impact, or experimental results required for acceptance at a conference. Instead, our review focuses on whether the writing contains misinterpretations of results, incorrect methodological descriptions, inaccurate citations, or hallucinations.
Therefore, we cross-check the manuscripts against the actual code and experimental results to precisely identify such issues.
The issues of our review are summarized as follows.
More detailed reviews for each paper are provided in the supplementary material.

Positive aspects are that none of these papers contained citations to non-existent works or developed invalid methods, such as test-data-leaking methods.
The issues in these papers are as follows:

\begin{tcolorbox}[colback=gray!10, colframe=gray!40!black, title=Issue1 :,breakable]
      Frequent Irrelevant Citations.
\end{tcolorbox}

We found that there are some irrelevant citations in these papers. This issue arises when adding new BibTeX entries that are not included in the baseline papers. (The reason for this is discussed in detail in Writing Risk 2 of \Cref{sec:risk}.)

\begin{tcolorbox}[colback=gray!10, colframe=gray!40!black, title=Issue2 :,breakable]
      Ambiguous Method Descriptions.
\end{tcolorbox}


We found that while the method descriptions are generally accurate, they still contain ambiguities.
For example, in the LoCoOp extended paper, the parameter appearing around Line 156 is not clearly explained, making it difficult to fully understand the method.
Similarly, in Min-K\%++ extended paper (Lines 126–129), although the corresponding code exists, the process is implemented as an optional component and is not actually utilized.  This occurs because the coding agent makes numerous modifications during Stage 2 in the Experimental Phase, increasing code complexity. This suggests that accurately transferring experimental code into a faithful methodological description remains an open challenge.

\begin{tcolorbox}[colback=gray!10, colframe=gray!40!black, title=Issue3 :,breakable]
      Misinterpretation of Figure Results.
\end{tcolorbox}

These papers include the overinterpretation of the figure results, making unsupported claims that appear plausible.
For instance, in Min-K\%++ extended paper (L177) and LoCoOp extended paper (L160–162), they report findings not evident from the figures. This highlights that precise result interpretation remains difficult for current AI Scientist systems.

\begin{tcolorbox}[colback=gray!10, colframe=gray!40!black, title=Issue4 :,breakable]
    Descriptions of Auxiliary Experiments That Were Never Conducted.
\end{tcolorbox}

We found several cases where these papers describe auxiliary experiments that were never actually conducted, such as in LoCoOp extended paper (Lines 183–184) and GL-MCM extended paper (Lines 208–213).
This issue occurred even though the writing agent was explicitly instructed not to include nonexistent experimental results.
This problem is especially tricky because hallucinations do not appear in the main results, which are easy to notice, but they often appear in parts like ablation or analysis. Therefore, even human reviewers might not notice them unless they carefully check the draft.
Such cases illustrate that the risk of hallucination remains inherent in the current system.

